% !TeX spellcheck = en_US
\documentclass[french]{yLectureNote}

\title{Électromagnétisme 1 PS}
\subtitle{Physique}
\author{Paulhenry Saux}
\date{\today}
\yLanguage{Français}
%lionels.calmels@cemes.fr
\professor{F.Pettinari}
\usepackage{graphicx}%----pour mettre des images
\usepackage[utf8]{inputenc}%---encodage
\usepackage{geometry}%---pour modifier les tailles et mettre a4paper
%\usepackage{awesomebox}%---pour les boites d'exercices, de pbq et de croquis ---d\'esactiv\'e pour les TP de PC
\usepackage{tikz}%---pour deiffner + d\'ependance de chemfig
% \usepackage{tabularx}%---pour dimensionner automatiquement les tableaux avec variable X
\usepackage{awesomebox}%---Pour les boites info, danger et autres
\usepackage{menukeys}%---Pour deiffner les touches de Calculatrice
\usepackage{fancyhdr}%---pour les en-t\^ete personnalis\'ees
\usepackage{blindtext}%---pour les liens
\usepackage{hyperref}%---pour les liens (\`a mettre en dernier)
\usepackage{caption}%---pour la francisation de la l\'egende table vers Tableau
\usepackage{pifont}
\usepackage{array}%---pour les tableaux
\usepackage{lipsum}
\usepackage{yFlatTable}
\usepackage{multicol}
\newcommand{\Lim}[1]{\lim\limits_{\substack{#1}}\:}
\renewcommand{\vec}{\overrightarrow}
\newcommand{\N}[0]{\mathbb{N}}
\newcommand{\dd}{\mathrm{d}}
\newcommand{\norm}[1]{||\vec{#1}||}
\newcommand{\fo}{\psi(\vec{r},t)}
\newcommand{\foe}{\psi(\vec{r},t)\*}
\newcommand{\HH}{\hat{H}}
\newcommand{\hb}{\hbar}
\newcommand{\lap}{\nabla^2}
\newcommand{\lapcc}{\frac{\partial^2 }{\partial x^2}+\frac{\partial^2 }{\partial y^2}+\frac{\partial^2 }{\partial z^2}}
\newcommand{\E}{\vec{E}(M)}
\setcounter{chapter}{1}
\begin{document}
\chapter{Distribution de charges}
\section{Répartition}
\begin{itemize}
 \item Atome : Noyau de diamètre \(10^{-15}\) + Électron \(10^{-10}\)
 \item Molécules : On pourra aussi avoir affaire à des distributions de charge qui ne sont uniforme (et qui dépendent par exemple du rayon )es entre atomes  :  \(10^{-10}\)\marginInfo{La matière peut \^etre globalement neutre mais pas localement. Il y aura toujours des champs électriques}.
\end{itemize}
\section{Les différentes distributions de charge}
\subsection{Densité de charge}
Si une charge macroscopique Q est répartie uniformément\marginWarning{On pourra aussi avoir affaire à des distributions de charge qui ne sont uniforme (et qui dépendent par exemple du rayon )} dans un volume V, alors on parle de distribution volumique de charge uniforme et la densité volumique vaut \(\rho_c = \frac{Q}{V}\). De la m\^eme façon, on a \(\sigma_c = \frac{Q}{S}\) ou \(\lambda_c = \frac{Q}{L}\).
\subsection{Symétrie d'une distribution de charge}
\subsubsection{Plan de symétrie (\(\pi^+\))}
\begin{itemize}
 \item Si \(\rho_c(x,y,z) = \rho_c(x,y,-z)\) alors le plan \((Oxy)\) est \(\pi^+\).
 \item le plan médiateur entre 2 particules de m\^eme charge.
 \item Cylindre uniformément chargé en volume : Plan médiateur, tout plan contenant l'axe de rotation
 \item Cylindre uniformément chargé en volume infiniment long : Plan médiateur, plans // au plan médiateur, tout plan contenant l'axe de rotation
 \item Sphère uniformément chargé en volume : Plans passant par le centre.
\end{itemize}
\subsection{Plans d'anti-symétrie (\(\pi^-\))}
\begin{itemize}
 \item Distribution volumique de charge tel que \(\rho_c(x,y,z) = -\rho_c(x,y,-z)\) alors le plan \(Oxy\) en est un
 \item le plan médiateur entre 2 particules de charge opposée.
 \item Une sphère coupée en 2 demie-sphères de charge opposée : plan coupant les 2 demie-sphères.
\end{itemize}
\subsection{Plans quelconques}
Ils présentent une symétrie géométrique mais pas de symétrie de charge
\subsection{Invariance d'une distribution de charge}
\subsubsection{Invariance par translation}
\begin{itemize}
 \item Si \(\rho_c(x,y,z) = \rho_c(x,y,z')\), la distribution est invariante par toute translation // à \((Oz)\).
 \item Cylindre infini d'axe Oz : Axe Oz
\end{itemize}
\subsubsection{Invariance par rotation}
\begin{itemize}
 \item Distribution de charge telle que \(\rho_c(\rho, \varphi, z) = \rho_c(\rho, \varphi', z)\) est invariant par rotation autour  de (Oz)
 \item Cylindre d'axe Oz : Invariance par rotation autour de \((Oz)\).
\end{itemize}
 \end{document}
